\section{Programs, attachments and acceptance}
\label{sec:attachment}

The growth studies keep elementary token operations fixed within each paired
comparison. Across studies, the supplied control language becomes richer:
first a sequence entering an old machine, then returning calls, nested
compositions, recursion, and chosen argument and return interfaces.

\subsection{The exact construction}

Programs are represented by finite directed multigraphs with observation and
action labels. Maps preserve endpoints and labels. We designate maps injective
on states and transitions as cofibrations. An attachment combines a retained
library $\mathcal L$ and a new program cell $B$ along their declared interface
$A$:
\[
\begin{array}{ccc}
 A & \hookrightarrow & B \\
 \big\downarrow & & \big\downarrow \\
 \mathcal L & \hookrightarrow & \mathcal L\amalg_A B .
\end{array}
\]
Both maps from $A$ are injective. The pushout takes the disjoint union and
identifies exactly the interface images. The implementation records the four
maps and content hashes, reconstructs the quotient on replay, and constructs
the unique mediating map for compatible maps out of $\mathcal L$ and $B$.
This is an ordinary graph pushout; no Quillen model structure is asserted.

The initial experiment permits only entry into old control, without a return.
Later boundaries include typed procedure entry and return ports. The retained
body is shared, not copied into each call. The interpreter decodes expressions
and call targets from the quotient graph. Old bodies and dependency hashes
cannot be changed. Several new helpers and a task entry procedure, called the
root, may be attached successively as one proposal.

A graph certificate establishes the attachment, not the new task's correctness.
For the later interface studies, a proposal must also typecheck, produce exact
outputs, return without error, consume all input, and preserve earlier tasks.
Search stops at its first exact training fit. Withheld validation then accepts
or rejects that proposal; a failure does not provide another training example.
In the cumulative A/B studies it ends the sequence, leaving later tasks blocked.
The whole accepted proposal is retained. Size does not rank successful proposals
or decide which helpers survive.

\subsection{Starting machines and permitted programs}

The returning-call study acquires five instruction fragments by bounded word
search: copy, skip and duplicate one token, swap a pair, and reverse a triple.
They use cursor advance, emission and two local registers. The nested-composition
study freezes those fragments and six whole-input procedures acquired in that
curriculum, then tests new arrangements of them.

The later interface studies begin with nine acquired fragments and three
\emph{supplied} recursive programs. Those programs have no arguments and differ
in their local token operations. They offer examples of related code from which
an abstraction can be drawn. Each proposer receives the same basis; learned
helpers are not exchanged between the independent arms or data seeds.

Helpers in these studies may have up to four Boolean or procedure arguments,
and return a Boolean or unit (no returned value). Procedure arguments name
procedures that take no arguments and return unit. Available constructs are
calls, recursion, sequencing,
conditionals, lexical bindings, Boolean operations and equality of the next
two tokens. There is no token literal, rewind or subtree instruction.
The root binds and composes helpers; input inspection and recursion belong
inside helpers. Proposals contain at most four definitions of at most 180
expression nodes each. Execution is limited to 60,000 steps and call depth 150,
with input progress required before recursive reentry. The finite program graph
therefore runs with a supplied call stack; it is not a finite-state machine
without auxiliary memory.

\subsection{The two proposal methods}

The mechanical algorithm changes with the language. The instruction comparison
uses breadth-first word search with output-prefix rejection and merging of
equal training configurations. The nested-composition comparison enumerates
normalised expression trees by call count and caches exact procedure results.
The first recursive comparison uses bounded counterexample-guided constraint
solving over control graphs.

For the later interface and control tasks, mechanical search first tries calls
to retained procedures. It then compares concrete source programs, replaces
differing calls to procedures without arguments with procedure parameters, and
checks that the abstraction can recover each source. Typed transformations add
Boolean arguments, return values and their uses. The expanded search also
inserts up to three effects before or after expressions, preserving returned
values through lexical bindings.

This search is not exhaustive over the accepted language. It rotates in batches
of sixteen among argument fillings and at most 512 edited templates. Unedited
templates receive at least one third of the batches and are never evicted;
edited templates can be evicted after 256 fillings. Candidates are screened
with a 12,000-step limit before the common 60,000-step execution check. The
recorded best training prefix does not guide enumeration, and parameterised
helpers are callable but are not all reused as new control templates. Private
reference programs confirm that each frontier task has a solution in this
mechanical grammar; they are withheld from both proposers.

The LLM is \texttt{gpt-5.6-sol}, with medium reasoning. It receives the language,
retained code, public training pairs and earlier public tasks, without task
names, target algorithms, future examples or desired signatures. Each retry
receives the last proposal and a compiler error or two public counterexamples.
Calls use the existing Codex access with tools disabled; the host accepts only
structured replies and does not repair them. The feedback is explicitly carried
between calls rather than relying on hidden session state.

Both methods work under a library-growth objective: try retained helpers,
factor common code, expose useful differences as arguments, and keep new glue
small without sacrificing a useful interface. Mechanical construction encodes
these preferences in search order; the LLM receives them as instructions.
They are preferences, not a hard size gate.

\subsection{What counts as reuse}

Direct reuse requires an unchanged helper hash, actual execution on the new
task, and a loss of correctness when its body is disabled. Recursive reuse
additionally requires recursive reentry in the execution trace. Merely keeping
a helper, or invoking only its leaf branch, is insufficient.

An interface extension is a different result. Here a new helper exposes choices
that the old one fixed internally. We check whether fixing the added arguments
recovers the old abstract syntax tree, including recursive bindings and effect
order. This post hoc check neither changes admission nor turns replacement of
an old cell into direct execution of it.

For the interface studies, the size diagnostic is
\begin{equation}
 \Delta d_k=\sum_{b\in\mathcal N_k} n(\operatorname{body}(b)),
 \label{eq:growth-size}
\end{equation}
where $\mathcal N_k$ contains every new helper and root accepted at task $k$,
and $n$ counts expression nodes. Parameter declarations are excluded. A drop
means less new code, not a smaller total library. This measure differs from
the rules-plus-instructions count in Eq.~\eqref{eq:fe} and the earlier modular
language's instruction, constructor and interface units.

\section{Where mechanical search was sufficient}
\label{sec:growth-ladder}

The earliest attachment pilot improved on cold evolutionary synthesis, but its
model arm allowed richer glue than its mechanical arm. It did not establish
an LLM advantage. The following comparisons matched the accepted language
and progressively broadened the compositions (Table~\ref{tab:growth-ladder}).

\begin{table}[htbp]
\centering
\small
\begin{tabular}{p{0.39\linewidth}rr}
\toprule
Experiment & Mechanical & LLM with harness \\
\midrule
Instruction glue: accepted conditions & 12/12 & 12/12 \\
\quad Total proposal time & 0.99 s & 125.75 s \\
Returning calls: accepted conditions & 24/24 & 24/24 \\
\quad Total proposal time & 20.63 s & 673.96 s \\
Nested composition: accepted tasks & 10/10 & 7/10 \\
\quad Total proposal time & 14.55 s & 674.79 s \\
Recursive control: acquisition & 95.88 s & 156.62 s \\
\quad Subsequent direct reuse & 0.28 s & 132.21 s \\
Chosen interfaces: accepted tasks & 6/6 & 6/6 \\
\quad Direct helper transfers & 3/3 & 3/3 \\
\bottomrule
\end{tabular}
\caption{Earlier comparisons. Counts are task/seed conditions, not independent
families. Times include proposal generation and training checks but exclude
later verification. Budgets differ between studies. In the chosen-interface
study, a shared mechanical search supplied every transfer binding, including
those in the LLM arm.}
\label{tab:growth-ladder}
\end{table}

Both methods passed all hidden cases in the matched instruction study and the
study of returning calls. The latter grew five fragments through twelve tasks
on two seeds. The final machine called fourteen procedures, including its root, with
seven procedure levels observed. Its seventeen stored entries used 98 term
units, versus 618 when each entry was separately inlined. Outputs and work
counts agreed with the inlined execution: sharing saved storage, not runtime.

Broader nested pipelines and branches still favoured enumeration. The hardest
tested composition took 4,053,996 candidates and 11.07 seconds. Three model
conditions timed out at 120 seconds; all seven accepted model programs passed
their withheld tests. Some long generating expressions also had shorter correct
solutions, so generating size alone was not a reliable difficulty measure.

Both methods learned a recursive traversal and rebound it on the next task,
but that study supplied the context bit and callback slots. An earlier pilot
had failed through mechanical memory limits and model timeouts. The successful
follow-up changed both the symbolic execution bound and the deadline, from
120 to 300 seconds.

The next study let proposers choose helper signatures in three families. Both
solved all six tasks. The common host search then found every transfer binding;
the model was responsible for its acquired helpers, not those later bindings.
A preceding trial with unequal library-growth objectives was discarded and is
not included in these results.

\subsection{Changes in recursive layout}

With generic effect insertion added to mechanical search, a 180-second pilot
tested four layouts. Both methods acquired the two easier layouts. Only the
model acquired one with token pairs before, between and after two recursive
children, in 65.34 seconds; a host-selected binding then transferred it.
Both failed the layout requiring information returned from a child. On a
fresh data seed neither method repeated the three-payload acquisition.

A fixed follow-up used seeds 2, 3 and 4, two model starts per seed, and one
deterministic mechanical run per seed. Each task allowed 600 seconds, six
model replies or eight million mechanical candidates. The model solved both
tasks in all six runs; mechanical search timed out on all three acquisitions.
Four model libraries supported direct recursive reuse through host-selected
bindings. The other two required interface extensions. No successful stage
needed a third reply, and five acquisitions fit the earlier 180-second limit.
The stronger evidence came mainly from repeated starts, not from showing that
more conversational turns were necessary. This still tested a model plus a
mechanical binder, motivating the independent proposal processes below.

\section{Independent mechanical and LLM proposals}
\label{sec:direct-ab}

In the direct A/B study each arm starts from its own copy of the original
basis and authors every attachment. The LLM gets no host-generated binding.
Both arms use the same accepted language, compiler, graph verifier and examples.
On seeds 31 and 32 they attempt six successive tasks: copy then duplicate leaves
in a layout with a middle payload; repeat for a layout with leading, middle
and trailing payloads; then repeat when a payload operation depends on a
child's returned information. A payload here is an ordinary token pair placed
beside the recursive children.

Each task has 26 training and 26 withheld validation examples, with 40 hidden
and thirteen stress cases evaluated after all choices freeze. Token pools are
disjoint; stress trees reach 128 leaves. The allowance is 300 seconds per task,
at most four model replies or eight million mechanical candidates. Failure
blocks later tasks in that arm without importing a solution.

\begin{table}[htbp]
\centering
\small
\begin{tabular}{lrrrr}
\toprule
Method & Accepted & Failed attempts & Blocked & Added nodes per seed \\
\midrule
Mechanical & 4 & 2 & 6 & $36,6$ \\
LLM & 8 & 2 & 2 & $27,6,35,6$ \\
\bottomrule
\end{tabular}
\caption{Direct A/B: twelve scheduled outcomes per arm. Both seeds have the
displayed size sequence. Blocked tasks are unattempted, not independent failures
or zero-cost solutions.}
\label{tab:direct-ab}
\end{table}

Mechanical search acquired the first parser in 36.1 and 38.0 seconds and its
transfer binding in about 0.28 seconds. The model took 67.6 and 109.3 seconds
for acquisition and 223.6 and 131.5 seconds for transfer. At task 3 the result
reversed: search timed out after 1,888,660 and 2,005,511 candidates, while the
model acquired the larger parser in 61.6 and 44.8 seconds, then supplied its
reuse calls in 13.6 and 12.4 seconds.

The two mechanical and four model transfers all used unchanged recursive
helpers on every hidden case. Traces confirm recursive reentry, and disabling
each helper reduced correctness from 40/40 to 0/40. Task 3 did, however, add
a new parser rather than reuse the task 1 parser. The repeated size drops
describe reuse within task pairs, not transfer across every change in format.

The model failed task 5 on both seeds. On seed 31 a structural shortcut passed
21/26 training examples, and the next reply timed out. On seed 32 a syntax
repair achieved 26/26 training but only 23/26 validation; all four old tasks
still passed. Its graph certificates were valid, but its new behaviour was
wrong. No withheld counterexample was returned to the proposer.

\section{A harder task with twenty-minute budgets}
\label{sec:hard-ab}

The next comparison combined three payload positions with information returned
from both children. An internal region consists of an equal-token marker pair,
a leading payload, a left child, a middle payload, a right child and a trailing
payload. Leaves and payloads are unequal pairs. The target swaps the leading
pair. It swaps the middle pair when the left child has an odd leaf count,
otherwise keeping only its first token. It keeps the trailing pair's first
token when the whole region has an odd leaf count, otherwise swapping it.
Task 1 copies leaves; task 2 duplicates their tokens. Proposers see examples,
not this rule or its parity interpretation.

Both arms start afresh on seeds 41 and 42. Each task has 44 training and 44
withheld validation examples, covering every binary shape through five leaves
plus combs and random trees with up to ten leaves. Forty hidden cases have 12--96
leaves; thirteen stress cases include 128-leaf trees. The deadline is 1,200
seconds per task, with ten model replies or 32 million mechanical candidates.
The model, prompt, search policy and acceptance rules are unchanged. The task,
examples and allowance all changed, so this is not an isolated test of extra
time on the preceding task.

\begin{table}[htbp]
\centering
\small
\begin{tabular}{rllrl}
\toprule
Seed & LLM acquisition & LLM task 2 & Added nodes & Mechanical acquisition \\
\midrule
41 & 709.9 s; 7 replies & 45.3 s; 1 reply & $77,84$ & 9,261,328; timeout \\
42 & 591.7 s; 5 replies & 70.9 s; 2 replies & $100,12$ & 10,377,878; timeout \\
\bottomrule
\end{tabular}
\caption{Hard A/B. Mechanical entries count synthesis candidates before the
twenty-minute limit. Both mechanical acquisitions failed, so neither dependent
task ran. All four model tasks passed acceptance and the later tests.}
\label{tab:hard-ab}
\end{table}

The model's successive acquisition scores were $2,2,6,2,6,37,44$ out of 44 on
seed 41, and a syntax error followed by $2,6,33,44$ on seed 42. Both winning
proposals exceeded the previous time and reply limits and passed all 44
validation cases.

\begin{samepage}
\paragraph{Seed 42: direct reuse.}
The acquisition introduced \texttt{Walk5Cut}, a recursive helper with procedure
arguments. The successful task 2 root was
\begin{verbatim}
(seq (call Walk5Cut (ref C02) (ref C02) (ref F004) (ref C03)) unit)
\end{verbatim}
\end{samepage}
Here \texttt{C02} duplicates a pair, \texttt{F004} swaps it and \texttt{C03}
keeps its first token. The helper's hash was unchanged; it recursed on all
forty hidden cases, and disabling it gave 0/40 correct results. The model
corrected an initially wrong binding after public feedback. Another acquired
helper, \texttt{FirstPair}, duplicated an existing fragment and was unused on
task 2. The $100\to12$ drop therefore shows reuse, not minimal acquisition.

\paragraph{Seed 41: interface extension.}
The first helper, \texttt{walk6}, had two Boolean arguments but fixed copying
at leaves. Its replacement, \texttt{walk6f}, adds a procedure argument
\texttt{emit}. Setting \texttt{emit} to the old copy fragment \texttt{C01}
and specialising the recursive calls recovers the old syntax tree exactly,
including lexical bindings and effect order. A separate auditor checked this
after selection; neither proposer nor the acceptance code used it. All 564
comparisons across four Boolean settings and 141 original-task inputs agreed
on output, cursor and returned value. Argument passing can still change the
interpreter step cost.

The old helper executed on zero task-2 hidden cases. This is verified interface
extension, not direct reuse or an extra pushout certificate. Its $77\to84$
size sequence has no immediate drop. The two-task protocol did not test whether
later calls would repay that extension.

\section{The later value of an extended interface}
\label{sec:transfer-value}

A separate retention assay followed all six model libraries from the
600-second repeated study, including its two extensions. It was not another
proposer A/B. Three unused leaf operations---swap, keep first and keep
second---were tested on data seeds 21, 22 and 23, without changing the recursive
format. The same mechanical adapter search ran on snapshots before and after
task 2. It could choose a retained procedure and bind its arguments, but could
not synthesise a new helper. Each search allowed ten seconds and 10,000
candidates. Nine fresh LLM controls started from the original basis, with
300 seconds and four replies per task.

Before their extensions, the two affected libraries solved 0/18 future tasks;
afterwards they solved 18/18. The failures exhausted the allowed adapters,
not every possible composition. Four libraries with already useful interfaces
solved 36/36 from either snapshot, and all nine fresh controls succeeded.
The extended helpers kept their hashes and were necessary in every corresponding
hidden case under body elision. All these transfer bindings came from search.

\begin{table}[htbp]
\centering
\small
\begin{tabular}{lrrrrr}
\toprule
Library history & Task 1 & Task 2 & Task 3 & Task 4 & Task 5 \\
\midrule
Extension to \texttt{P6} & 31 & 45 & 7 & 7 & 7 \\
Extension to \texttt{L1} & 33 & 47 & 8 & 8 & 8 \\
Direct reuse (two histories) & 42 & 8 & 8 & 8 & 8 \\
Direct reuse & 45 & 7 & 7 & 7 & 7 \\
Direct reuse & 35 & 6 & 6 & 6 & 6 \\
\bottomrule
\end{tabular}
\caption{New expression nodes per task. Tasks 1--2 are the six frozen
histories; tasks 3--5 were chosen before the follow-up ran. Counts agree across
the three new data seeds, which do not represent new interface discoveries.}
\label{tab:extension-cost}
\end{table}

Fresh solutions used 24--61 nodes, versus 6--8 for adapters. Charging the
entire task-2 addition, including its root, later savings covered the 45-node
\texttt{P6} extension after one, two and two tasks on the three seeds. The
47-node \texttt{L1} extension was covered after one, three and two. This payback
is relative to the observed fresh solutions, not a minimum description bound.
The before/after comparison gives the size curves a behavioural interpretation:
the added interface made those cheap calls possible. The evidence concerns
new leaf operations in one format, not repeated discovery across unrelated tasks.

\section{Verification and reproducibility}
\label{sec:growth-records}

For the two direct comparisons, all scheduled outcomes were frozen before
hidden evaluation. Fresh processes checked proposal provenance, separate
library histories, graph certificates and failed prefixes. Each accepted
program also ran on all binary shapes through seven leaves with two token
assignments, and on 81 pairs of concatenated regions to test return at an
internal boundary. Table~\ref{tab:growth-audit} gives the recorded totals.

\begin{table}[htbp]
\centering
\small
\begin{tabular}{lrrr}
\toprule
Check & Direct A/B & Hard A/B & Retention assay \\
\midrule
Accepted task instances & 12 & 4 & 99 \\
Hidden cases correct & 480/480 & 160/160 & 3,960/3,960 \\
Stress cases correct & 156/156 & 52/52 & 1,287/1,287 \\
Shape checks correct & 4,728/4,728 & 1,576/1,576 & 39,006/39,006 \\
Continuation checks correct & 972/972 & 324/324 & 8,019/8,019 \\
Accepted attachment certificates & 18 & 8 & 109 \\
\bottomrule
\end{tabular}
\caption{Finite verification of accepted programs, not independent discovery
trials or a proof for all inputs. The retention assay also records eighteen
adapter failures. Two further graph certificates in the direct A/B belonged
to a behaviourally rejected proposal and are excluded from its accepted count.}
\label{tab:growth-audit}
\end{table}

Total proposal time was 675.1 seconds mechanically and 1,250.8 seconds for
the model in the direct A/B; the hard A/B used 2,400.6 and 1,417.8 seconds.
The latter tested 19,639,206 synthesis candidates and completed fifteen model
replies. Neither hard-trial arm hit its sampled 640 MiB local process-group
limit. These resource measures cover different amounts of completed work and
do not include remote model hardware. Detailed usage and earlier study budgets
are given in the individual reports.

The unified reproduction package is described in
\path{transduction/REPRODUCIBILITY.md}; the study-to-evidence map is in
\path{transduction/manuscript/README.md}. It includes the register rerun and
the later studies, with sources, data, proposals, receipts, program histories,
certificates and replay commands. Historical records are preserved unchanged.
Two early auxiliary audits whose source copies were missing were regenerated;
every previously reported result was required to match. Their new reports are
separate from the originals. Replay checks the recorded source versions and
runs without model access. The package is a local deliverable, not yet a public
archival deposit.
